%───────────────────────────────────────────────────────────────────────────────
% ESTILO PARA ALGORITMOS CON COLORES TEMÁTICOS CATPPUCCIN 
%───────────────────────────────────────────────────────────────────────────────

% Definir comandos de estilo que toman un argumento
\newcommand{\algkw}[1]{\textcolor{keywordcolor}{#1}}      % palabras clave (antes CtpMauve)
\newcommand{\algarg}[1]{\textcolor{maintext}{#1}}          % argumentos (antes CtpText)
\newcommand{\algcomment}[1]{\textcolor{commentcolor}{#1}}  % comentarios (antes CtpOverlay2)
\newcommand{\algnum}[1]{\textcolor{operatorcolor}{#1}}     % números de línea (antes CtpSky)
\newcommand{\algcap}[1]{\textcolor{linkcolor}{#1}}         % caption (antes CtpBlue)

% Aplicar los estilos a algorithm2e (¡sin la barra invertida!)
\SetKwSty{algkw}
\SetArgSty{algarg}
\SetCommentSty{algcomment}
\SetNlSty{algnum}{}{}
\SetAlCapSty{algcap}

% Caja decorativa (opcional)
\tcolorboxenvironment{algorithm}{
  colback=background!0,          % fondo transparente (antes CtpBase!0)
  colframe=surface0,             % marco sutil (antes CtpSurface0)
  arc=0pt,
  boxrule=0pt,
  left=6pt,
  right=6pt,
  top=6pt,
  bottom=6pt,
  before skip=10pt,
  after skip=10pt,
  drop fuzzy shadow=black!5
}

% ── Color y grosor de líneas de algoritmos ────────────────────────────────────
\setlength{\algoheightrule}{0.6pt}       % línea superior e inferior
\setlength{\algotitleheightrule}{0.5pt}  % línea bajo el caption

\makeatletter
% Usamos operatorcolor (CtpSky) para las líneas, o puedes definir uno específico
\newcommand{\algolinecolor}{operatorcolor}    % antes CtpSky

\def\@algocf@pre@ruled{%
  {\color{\algolinecolor}\hrule width\algocf@ruledwidth
   height\algoheightrule depth0pt}%
  \kern\interspacetitleruled}

\def\@algocf@post@ruled{%
  \kern\interspacealgoruled
  {\color{\algolinecolor}\hrule width\algocf@ruledwidth
   height\algoheightrule}}

\renewcommand{\algocf@caption@ruled}{%
  \box\algocf@capbox
  \kern\interspacetitleruled
  {\color{\algolinecolor}\hrule width\algocf@ruledwidth
   height\algotitleheightrule depth0pt}%
  \kern\interspacealgoruled}
\makeatother

% Línea vertical izquierda (para entorno ruled con vlined)
\makeatletter
\renewcommand{\algocf@Vline}[1]{%
  \strut\par\nointerlineskip%
  \algocf@push{\skiprule}%
  \hbox{{\color{\algolinecolor}\vrule}%
    \vtop{\algocf@push{\skiptext}%
      \vtop{\algocf@addskiptotal #1}\Hlne}}\vskip\skiphlne%
  \algocf@pop{\skiprule}%
  \nointerlineskip}
\makeatother

\makeatletter
\renewcommand{\algocf@Vsline}[1]{%
  \strut\par\nointerlineskip%
  \algocf@bblockcode%
  \algocf@push{\skiprule}%
  \hbox{{\color{\algolinecolor}\vrule}%
    \vtop{\algocf@push{\skiptext}%
      \vtop{\algocf@addskiptotal #1}}}%
  \algocf@pop{\skiprule}%
  \algocf@eblockcode%
}
\makeatother